% Source: https://artofproblemsolving.com/community/c3h1870491_inscribed_semicircle_radius
% Stable identity: aops-topic-1870491-q01
(a) Let $O$ be the centre of the semicircle; it is the midpoint of the diameter, so it lies on $BC$. Let the arc touch $AB$ at $T_1$ and $AC$ at $T_2$. A radius drawn to a point of tangency is perpendicular to the tangent line, so $OT_1\perp AB$ and $OT_2\perp AC$ with $OT_1=OT_2=r$. In particular $O$ is equidistant from lines $AB$ and $AC$, so it is the point where the bisector of angle $A$ meets $BC$.

\begin{asy}
usepackage("amsmath");
size(7.5cm);
defaultpen(fontsize(10pt));
pair B=(0,0),C=(14,0),A=(5,12),O=(6.5,0);
pair T1=B+dot(O-B,unit(A-B))*unit(A-B), T2=C+dot(O-C,unit(A-C))*unit(A-C);
void ra(pair P,pair Q,pair R,real s){pair u=unit(P-Q)*s,v=unit(R-Q)*s;draw(Q+u--Q+u+v--Q+v,linewidth(0.6));}
fill(arc(O,6,0,180)--cycle,rgb(0.88,0.92,1));
draw(arc(O,6,0,180),blue+linewidth(1));draw((O.x-6,0)--(O.x+6,0),blue+linewidth(1));
draw(A--B--C--cycle,linewidth(1));
draw(A--O,dashed);draw(O--T1,dashed);draw(O--T2,dashed);
ra(O,T1,B,0.6);ra(O,T2,C,0.6);
dot(O);dot(T1);dot(T2);
label("$A$",A,N);label("$B$",B,SW);label("$C$",C,SE);label("$O$",O,S);
label("$13$",(A+B)/2+(-0.4,0.6),NW);label("$15$",(A+C)/2,NE);label("$14$",(10,0),S);
label("$r$",(O+T1)/2+0.45*unit(rotate(90)*(T1-O)),fontsize(9pt));label("$r$",(O+T2)/2+0.45*unit(rotate(-90)*(T2-O)),fontsize(9pt));
label("\mbox{}",(min(currentpicture,user=true)+max(currentpicture,user=true))/2,black); // keeps the text after the figure black when LaTeX typesets the labels inline
\end{asy}

The segment $AO$ splits triangle $ABC$ into triangles $ABO$ and $ACO$. Taking $AB$ and $AC$ as bases, their heights are $OT_1=r$ and $OT_2=r$. Hence
\[
K=[ABO]+[ACO]=\dfrac{1}{2}\,AB\cdot r+\dfrac{1}{2}\,AC\cdot r=\dfrac {r}{2}\,(AB+AC),
\qquad\text{so}\qquad r=\dfrac{2K}{AB+AC}.
\]

(b) By Heron's formula with semiperimeter $s=\dfrac{13+14+15}{2}=21$,
\[
K=\sqrt{21\cdot(21-13)(21-14)(21-15)}=\sqrt{21\cdot8\cdot7\cdot6}=\sqrt{7056}=84 .
\]
So $r=\dfrac{2\cdot84}{13+15}=\dfrac{168}{28}=\boxed{6}$. (The centre is at distance $\dfrac{13}{28}\cdot14=6.5$ from $B$, and the diameter, from $0.5$ to $12.5$, fits inside $BC$.)
